\begin{proposition}[F4: a quantitative sufficient criterion for prosolvable quotients; Proved]
Let $G$ be a prosolvable profinite group. Suppose a real constant $C\geq0$
satisfies
\[
 \dim_{\mathbf F_\ell}H^1_{\mathrm{cont}}(G,M)
 \leq C\dim_{\mathbf F_\ell}M
\]
for every prime $\ell$ and every finite simple continuous
$\mathbf F_\ell[G]$-module $M$. Then $G$ has at most
$\lceil C\rceil+1$ topological generators. Conversely, $d$ generators
imply this bound with $C=d$.
\end{proposition}
\begin{proof}
Only sufficiency is new here. Put $D=\lceil C\rceil+1$.
For every finite quotient $F$ of $G$, inflation injects
$H^1(F,M)$ into $H^1(G,M)$ for an $F$-module $M$.
We show by induction on $|F|$ that every finite solvable quotient
$F$ is generated by $D$ elements, allowing identities in the tuple.
Choose a minimal nontrivial normal subgroup $M$ of $F$.
Solvability and minimality imply that $M$ is an elementary abelian
$\ell$-group and a simple module for $H=F/M$ under conjugation:
its derived subgroup, primary parts, and subgroup $\ell M$ are
characteristic, so minimality successively forces these assertions.
By induction choose a generating $D$-tuple of $H$.

Any lift of this tuple generates a subgroup $J$ with $JM=F$.
The intersection $J\cap M$ is an $H$-submodule, hence is either
$0$ or $M$. In the latter case $J=F$; in the former case $J$ is
a complement to $M$. Thus a nonsplit extension has a generating
lift immediately. In the split case identify $F=M\rtimes H$.
The complements are precisely the graphs of the crossed
homomorphisms $H\to M$. For the chosen tuple each complement
contains exactly one lifting tuple, so the number of nongenerating
lifts is $|Z^1(H,M)|$.

Restriction to $M$ and the splitting give, directly on cocycles,
\[
 H^1(F,M)\simeq H^1(H,M)\oplus\operatorname{End}_H(M).
\]
Indeed a cocycle vanishing on $H$ is exactly an $H$-equivariant
additive map $M\to M$, and a coboundary restricts to zero on $M$.
The endomorphism group contains its nonzero identity, so has
cardinality at least $\ell$. Consequently
\[
 |Z^1(H,M)|
 =|H^1(H,M)|\frac{|M|}{|M^H|}
 \leq \frac{|M|^{C+1}}{\ell}
 <|M|^D.
\]
There are $|M|^D$ lifting tuples, so a generating one exists.
This proves the finite induction.

For each open normal subgroup $N$ of $G$, the tuples in $G^D$
generating $G/N$ form a nonempty closed subset. Any finite
intersection contains the subset corresponding to the
intersection of the relevant normal subgroups, which is nonempty
by the finite result. Compactness therefore supplies one tuple
generating every finite quotient; its closed generated subgroup
is $G$. Necessity was proved in F3b by evaluating cocycles on a
generating tuple.
\end{proof}

\paragraph{The resulting geometric target.}
For the maximal prosolvable quotient of the local fundamental group,
it would suffice to prove a constant $C_A$, independent of the
finite simple local system $\mathcal M$ and its rank, with
\[
 \dim_{\mathbf F_\ell}H^1_{\mathrm{et}}
       (\Spec(A)\setminus\{\mathfrak m\},\mathcal M)
 \leq C_A\operatorname{rank}(\mathcal M).
\]
One only needs local systems whose monodromy is solvable; inflation
then bounds the cohomology of the prosolvable quotient. The argument
does not provide this uniform estimate. In particular, separate
finiteness on each finite cover has no uniform constant, as F3
demonstrates. Nonabelian simple chief factors of the full fundamental
group require an additional argument even after this target is met.
